% ----------------------------------------------------------------------
\section{Métricas (KPIs) y fórmulas}
\label{sec:kpis}

Para evaluar el desempeño y la aceptación del sistema, se han definido indicadores clave de rendimiento (KPIs). Estos métricas permiten cuantificar la eficiencia de los algoritmos de generación, la calidad de los horarios resultantes y la fiabilidad de las funciones auxiliares.

\begin{enumerate}
    \item TTS (Time-To-Schedule): Mide la latencia percibida en la generación.
    \[ TTS = t_{aplicado} - t_{click\_generar} \]
    El objetivo es mantener este valor por debajo de 3 segundos para $N=10$ combinaciones, asegurando una experiencia de usuario fluida durante la exploración de alternativas.

    \item Gaps/día: Evalúa la compacidad del horario generado.
    \[ Gaps_{prom} = \frac{\sum_{d} \text{minutos vacíos}(d)}{\text{\#días activos}} \]
    Se busca una tendencia a la baja en esta métrica, priorizando horarios con bloques de clases continuos para optimizar el tiempo del estudiante en la universidad.

    \item On-time reminders: Mide la fiabilidad de las notificaciones.
    \[ \text{Reliability} = \frac{\text{\#notifs abiertas } \ge X \text{ min antes}}{\text{\#notifs totales}} \times 100\% \]
    El objetivo es superar el 95\% de entregas puntuales, validando que el mecanismo de alarmas del sistema operativo no sea interrumpido por optimizaciones de batería.

    \item Adopción JSON: Indicador de viralidad y utilidad social.
    Se mide como el porcentaje de operaciones de exportación e importación exitosas respecto al total de intentos. Se espera un incremento sostenido durante las semanas de inscripción, validado mediante pruebas de extremo a extremo (E2E).
\end{enumerate}
